\documentclass[../main.tex]{subfiles}

\begin{document}

%% ======================================================================
%%  Section 5 – HDF5 Wrappers and Language Bindings
%% ======================================================================
\section{HDF5 Wrappers and Language Bindings}
\label{sec:wrappers}

\subsection{Boundary Model}

The SSP SIG's HDF5 Python Wrapper (H5PW) model~\cite{hdf5-python-wrapper-model}
defines ten mechanism families that generalize to any binding between HDF5 and
a managed runtime: \textbf{W1} native parser or codec compromise, \textbf{W2}
code in data, \textbf{W3} dynamic-loader boundary failure, \textbf{W4}
concurrency or lifecycle mismatch, \textbf{W5} resource exhaustion or
amplification, \textbf{W6} ABI,
dependency, or packaging drift, \textbf{W7} metadata, log, or artifact
exposure, \textbf{W8} trust-boundary confusion, \textbf{W9} misleading safety
signal, and \textbf{W10} external-reference traversal. The labels classify
mechanisms and carry no risk rating.

Bindings for C\#, Java, Python, JavaScript, Rust, and other languages mediate
HDF5's C interfaces into different object, memory, exception, packaging, and
deployment models. The Python cases below are boundary references, not an
exhaustive binding review; none establishes that a downstream framework
failure was caused by the binding through which HDF5 was reached.

\subsection{Python-Facing Cases}

All four advisories reached users through Python-facing workflows, but the
binding performed as specified and the missing decision was above it.

\begingroup
\small
\setlength{\tabcolsep}{3pt}
\renewcommand{\arraystretch}{1.12}
\begin{longtable}{@{}p{0.18\textwidth}p{0.61\textwidth}p{0.16\textwidth}@{}}
\hline
\textbf{Case} & \textbf{Boundary conclusion} & \textbf{Routing / family} \\
\hline
\endfirsthead
\hline
\textbf{Case} & \textbf{Boundary conclusion} & \textbf{Routing / family} \\
\hline
\endhead
\hline
\endfoot

CVE-2025-9905~\cite{cve-2025-9905} &
Keras legacy \code{.h5} loading executed serialized \code{Lambda} code despite
\code{safe_mode=True}. Application-level framework-object deserialization,
not HDF5 decode, was unsafe; the safety label did not govern this loader path. &
\code{APP-SEC-002}; W2, W9 \\
\hline

CVE-2026-1669~\cite{cve-2026-1669} &
Model loading followed a documented HDF5 external raw-storage reference to a
local path and disclosed its contents. The consumer lacked a policy deciding
whether this file could exercise that capability. &
\code{APP-SEC-003}; W10 \\
\hline

CVE-2026-0897~\cite{cve-2026-0897} &
Weight loading allocated from a file-declared shape without a budget. HDF5 can
legitimately describe a logical shape much larger than the useful payload, so
the consumer must enforce the budget. &
\code{APP-SEC-004}; W5 \\
\hline

CVE-2026-2492~\cite{cve-2026-2492} &
TensorFlow's integration searched an unsecured location for plugins. Dynamic
loading is an HDF5 mechanism; the defect was the embedding application's
search-path policy. &
\code{APP-SEC-005}; W3 \\
\hline
\end{longtable}
\endgroup

Together the cases instantiate H5PW \textbf{W8}: a binding faithfully exposes
an HDF5 capability, while the security decision---whether this process, with
this authority, should honor it for this file---belongs above the binding.
All five application-boundary findings in this report map to named W-families;
the vocabulary anticipated the boundary independently of this assessment.
Accordingly, \code{APP-S1} asks bindings to surface controls through their APIs
rather than leave them reachable only through environment variables.

\subsection{Coverage and Mitigation Routing}

No Java-specific adverse-event scenario, or additional binding-specific
scenario, was sufficiently developed for a formal finding. Java remains in the
ecosystem census and control-exposure scope; W1, W3, W4, and W6 are especially
applicable to its JNI boundary and provide a vocabulary for the next review.

The four cases remain formal application findings in
Appendix~\ref{app:findings-register}, not h5py or Python-binding findings.
Wrappers therefore receive no separate report-wide urgency: they participate
in the Phase~1 version, linkage, feature, and authority census and expose the
controls required by the applicable application package. Where a binding is
part of a deployment matching a formal finding, that finding's response clock
controls; mentioning a CVE here does not create a second priority.

\end{document}
